% =============================================================
% 4.3 การออกแบบองค์ประกอบในเกม
% โครงงาน: การออกแบบและพัฒนาเกมเดินเรือด้วยระบบสั่งการด้วยเสียงเพื่อส่งเสริมทักษะการสื่อสาร
% =============================================================

\section{การออกแบบองค์ประกอบในเกม}

องค์ประกอบภายในเกมได้รับการออกแบบเพื่อสร้างสภาพแวดล้อมจำลองที่สมจริง (Immersive Environment) และเอื้อต่อการเรียนรู้และการฝึกทักษะการสื่อสารภาษาอังกฤษ โดยแบ่งองค์ประกอบออกเป็น 3 ส่วนหลัก ดังนี้

\subsection{การออกแบบตัวละครและบทบาท (Characters and Roles)}

เกมจำลองสถานการณ์การทำงานร่วมกันบนห้องบังคับการเรือ โดยแบ่งบทบาทตัวละครออกเป็น 3 ตำแหน่งหลัก ดังนี้

\begin{enumerate}[label=\arabic*)]
    \item \textbf{กัปตัน (Captain):} ตัวแทนของผู้เล่น (Player Character) เป็นผู้มีอำนาจตัดสินใจสูงสุดและเป็นผู้ออกคำสั่งเสียงภาษาอังกฤษเพียงผู้เดียว
    \item \textbf{พลต้นหน (Navigator - AI):} ลูกเรือปัญญาประดิษฐ์ ทำหน้าที่ให้ข้อมูลด้านเส้นทาง เตือนภัยจากเรดาร์ และประเมินระยะห่างของสิ่งกีดขวาง เพื่อให้กัปตันนำข้อมูลไปตัดสินใจ
    \item \textbf{พลถือท้าย (Helmsman - AI):} ลูกเรือปัญญาประดิษฐ์ ทำหน้าที่รับคำสั่งจากกัปตันโดยตรงเพื่อบังคับพวงมาลัยหางเสือและคันเร่งเครื่องยนต์ พร้อมทั้งให้เสียงทวนคำสั่ง (Voice Feedback) เพื่อยืนยันความเข้าใจ
\end{enumerate}

\subsection{การออกแบบอุปสรรคและสิ่งกีดขวาง (Obstacles and Hazards)}

เพื่อให้เกมมีความท้าทายและป้องกันไม่ให้ผู้เล่นใช้การท่องจำเส้นทาง (Memorization) ระบบจึงใช้การสุ่มสร้างสิ่งกีดขวาง (Obstacle Generation System) รูปแบบของอุปสรรคที่ออกแบบไว้ ได้แก่

\begin{itemize}
    \item \textbf{สิ่งกีดขวางทางธรรมชาติ:} เช่น โขดหินใต้น้ำ เกาะ หรือภูเขาน้ำแข็ง ที่บีบบังคับให้ผู้เล่นต้องประเมินระยะห่างและสั่งการเลี้ยวหลบให้ทันเวลา
    \item \textbf{สิ่งกีดขวางแบบเคลื่อนที่:} เช่น เรือลำอื่นที่สัญจรไปมาในช่องแคบที่พลุกพล่าน ทำให้ผู้เล่นต้องฝึกไหวพริบและแก้ไขสถานการณ์เฉพาะหน้า
\end{itemize}

\subsection{การออกแบบกราฟิกและเสียงประกอบ (Graphic \& Audio Design)}

การออกแบบด้านทัศนศิลป์และเสียงมุ่งเน้นไปที่การสร้างความรู้สึกดื่มด่ำ (Immersion) ให้ผู้เล่นรู้สึกเหมือนกำลังบัญชาการเรือจริง ๆ

\begin{itemize}
    \item \textbf{กราฟิก (Graphic):} นำเสนอในมุมมองการเล่นแบบไอโซเมตริก (Isometric View) เพื่อให้ผู้เล่นสามารถมองเห็นตัวเรือ สภาพแวดล้อมทางทะเล และตำแหน่งของสิ่งกีดขวางโดยรอบได้อย่างชัดเจนในมุมกว้าง
    \item \textbf{เสียงแวดล้อม (Ambient Audio):} มีการใส่เสียงคลื่น เสียงลมทะเล และเสียงการทำงานของเครื่องยนต์เรืออย่างต่อเนื่อง เพื่อสร้างบรรยากาศการเดินเรือที่สมจริง
    \item \textbf{เสียงเอฟเฟกต์ (SFX):} ใช้เสียงสัญญาณเตือนภัย (Alarm) จากเรดาร์เมื่อเรือเคลื่อนเข้าใกล้สิ่งกีดขวาง เพื่อกระตุ้นความตื่นตัวและสร้างความกดดันด้านเวลาให้ผู้เล่นต้องรีบออกคำสั่ง
\end{itemize}